% 地球（综述）
% license CCBYSA3
% type Wiki

本文根据 CC-BY-SA 协议转载翻译自维基百科\href{https://en.wikipedia.org/wiki/Earth}{相关文章}。

\begin{figure}[ht]
\centering
\includegraphics[width=6cm]{./figures/5280dfaa768e0550.png}
\caption{“蓝色弹珠”，阿波罗 17 号，1972 年 12 月} \label{fig_Earth_1}
\end{figure}
地球是距太阳的第三颗行星，也是唯一已知能够孕育生命的天体。这得益于地球是一颗海洋世界——在太阳系中，唯一能够维持液态地表水的星球。地球上几乎所有的水都存在于全球性海洋中，覆盖了地球地壳的 70.8\%。剩余的 29.2\% 地壳为陆地，其中大部分集中在地球陆半球的大陆地块上。地球的大多数陆地区域至少具有一定湿度并被植被覆盖，而地球极地沙漠中的大型冰盖所储存的水量超过了地球地下水、湖泊、河流和大气水的总和。地球的地壳由缓慢移动的构造板块组成，彼此相互作用形成山脉、火山与地震。地球拥有液态外核，可产生磁层，能够偏转大部分具有破坏性的太阳风和宇宙射线。

地球拥有一个动态的大气层，这维持了地球的表面环境，并在陨石体及紫外线入射时提供保护。大气层主要由氮气和氧气组成。水汽在大气层中广泛存在，形成覆盖地球大部分区域的云层。水汽是一种温室气体，并与大气中的其他温室气体（特别是二氧化碳 $CO_2$）一道，通过捕获来自太阳光的能量，创造出液态表面水和水汽能够持续存在的条件。这一过程维持了目前 14.76 °C（58.57 °F）的平均地表温度，在正常大气压下，水在该温度呈液态。不同地理区域所捕获能量的差异（如赤道地区获得的阳光多于极地地区）驱动了大气和海洋环流，形成具有不同气候区的全球气候系统，并产生降水等多种天气现象，使碳和氮等组成成分得以循环。

地球呈近似椭球形，周长约为 40,000 公里（24,900 英里）。它是太阳系中密度最高的行星。在四颗类地行星中，地球最大且质量最高。地球距太阳约八光分（1 天文单位），沿轨道绕太阳运行一周约需一年（约 365.25 天）。地球自转一周所需时间略少于一天（约 23 小时 56 分钟）。地球自转轴相对于其绕太阳轨道平面的垂线有倾斜，从而产生四季更替。地球拥有一颗永久天然卫星——月球。月球以 384,400 公里（238,855 英里，约 1.28 光秒）的距离绕地球运行，其直径约为地球的四分之一。月球的引力有助于稳定地球自转轴，产生潮汐，并逐渐减缓地球自转。同样，地球的引力已使月球自转潮汐锁定，月球永远以同一面朝向地球。

地球与太阳系中大多数其他天体一样，约在 45 亿年前由早期太阳系中的气体和尘埃形成。在地球历史的最初十亿年中，海洋形成，随后生命在海洋中出现。生命在全球扩散，并持续改变地球的大气和地表，导致约 20 亿年前发生“大氧化事件”。人类约 30 万年前在非洲出现，随后扩散到地球的每一块大陆。人类依赖地球的生物圈和自然资源而生存，但对地球环境的影响却日益增加。如今，人类对地球气候与生物圈的影响已不可持续，威胁着人类自身及许多其他生命形式的生计，并导致物种广泛灭绝。

地球呈椭球形，周长约 40,000 公里（24,900 英里）。它是太阳系中密度最高的行星。在四颗岩质行星中，地球最大且质量最高。地球距离太阳约八光分（1 个天文单位），绕太阳公转一周约需一年（约 365.25 天）。地球绕自身轴心自转一周时间略少于一天（约 23 小时 56 分钟）。地球自转轴相对于其绕太阳轨道平面垂线存在倾斜，从而产生四季变化。地球拥有一颗永久天然卫星——月球，月球以 384,400 公里（238,855 英里，即 1.28 光秒）的距离绕地运行，其直径约为地球的四分之一。月球的引力有助于稳定地球自转轴、引发潮汐，并逐渐减缓地球自转。同样，地球的引力已使月球自转产生潮汐锁定，月球始终以同一面朝向地球。

地球与太阳系中大多数其他天体一样，约在 45 亿年前由早期太阳系中的气体和尘埃形成。在地球历史最初的十亿年中，海洋开始形成，随后生命在其中诞生。生命在全球扩散，并持续改变地球的大气层和地表，导致约二十亿年前发生“大氧化事件”。人类约在 30 万年前出现在非洲，并扩散至地球的每一块大陆。人类依赖地球的生物圈和自然资源生存，但对地球环境的影响日益加剧。目前人类对地球气候和生物圈的影响已不可持续，威胁着人类和许多其他生命形式的生存，并正在引发广泛的物种灭绝。
\subsection{词源}
现代英语单词 earth 经由中古英语发展而来，源自古英语名词 eorðe $^{[22]}$。在所有日耳曼语言中都有其同源词，由此重建出原始日耳曼语形式 erþō。在最早的文献中，单词 eorðe 用来翻译拉丁语 terra 和希腊语 gē 的多重含义：地面、泥土、陆地、人类世界、世界表面（包括海洋）以及整个地球本身。与罗马神话中的 Terra（或 Tellus）和希腊神话中的 Gaia 类似，“地球”可能在日耳曼多神教中也被视为一位拟人化的女神：晚期北欧神话中出现过 Jörð（“大地”），这位女巨人常被认为是雷神托尔的母亲 $^{[23]}$。

历史上，“earth”通常以小写书写。在中古英语早期，为表达“地球”这一明确含义，开始使用短语 the earth。到了早期现代英语时期，名词大写逐渐普及，the earth 也写作 the Earth，尤其是在与其他天体并列提及时。近来，这一名称有时直接写作 Earth，以与其它行星名称类比，然而 earth 和 the earth 的用法仍然常见 $^{[22]}$。不同出版风格各有规定：牛津拼写认为小写形式最为常见，而大写形式则为可接受的变体。另一种惯例是在作为名称使用时采用大写（如 “Earth's atmosphere”），但在前有定冠词时使用小写（如 “the atmosphere of the earth”）。在口语表达中几乎总是写作小写，例如“what on earth are you doing?” $^{[24]}$

名称 Terra（/ˈtɛrə/）有时用于科学写作；它也常见于科幻作品中，用来将人类居住的行星与其他星球区分开来 $^{[25]}$。而在诗歌中，Tellus（/ˈtɛləs/）则被用来表示地球的拟人形象 $^{[26]}$。Terra 也是一些拉丁语系语言（如意大利语和葡萄牙语）中“地球”的名称；在其他罗曼语中，该词的拼写略有变化，如西班牙语 Tierra 与法语 Terre。希腊诗意名称 Gaia 的拉丁化形式 Gaea 在英语中较为少见，尽管由于“盖亚假说（Gaia hypothesis）”，另一种拼写 Gaia 已变得常用，并读作 /ˈɡaɪ.ə/，而不是更传统的英语发音 /ˈɡeɪ.ə/ $^{[27]}$。

关于地球还有许多形容词。earthly（地上的、尘世的）源自 Earth。由拉丁语 Terra 衍生出 terran（/ˈtɛrən/）、terrestrial（/təˈrɛstriəl/），以及（经由法语） terrene（/təˈriːn/）；由拉丁语 Tellus 衍生出 tellurian（/tɛˈlʊəriən/）和 telluric $^{[28][29][30][31][32]}$。
\subsection{自然历史}
\subsubsection{形成}
\begin{figure}[ht]
\centering
\includegraphics[width=8cm]{./figures/409cc725593ef0b7.png}
\caption{对形成地球及其他太阳系天体的早期太阳系原行星盘的描绘} \label{fig_Earth_2}
\end{figure}
在太阳系中发现的最古老物质被测定为距今 $4.5682^{+0.0002}_{-0.0004}
$ 十亿年前$^{[33]}$。在距今 $4.54\pm0.04$ 十亿年前，原始地球已经形成$^{[34]}$。太阳系中的各个天体与太阳一起形成并演化。理论上，太阳星云通过引力坍缩从分子云中分割出一块空间，开始旋转并扁平化成一个环星盘，然后行星在该盘内与太阳一起生长。星云包含气体、冰粒和尘埃（包括原始核素）。根据星云理论，微行星通过吸积而形成，据估计原始地球可能需要约 7000 万年至 1 亿年才能形成$^{[35]}$。

对月球年龄的估计从距今 45 亿年到显著更年轻不等$^{[36]}$。主流假说认为，月球由一次撞击后从地球上抛出的物质吸积而成——一个约为地球质量 10\%、大小与火星相似的天体（命名为 忒伊亚 Theia）与地球发生碰撞$^{[37]}$。它以掠擦方式撞击地球，其部分质量与地球合并$^{[38][39]}$。大约在距今 40 亿年至 38 亿年前，晚期重轰炸**期间大量小行星撞击导致月球大范围表面环境发生显著改变，并据此推断地球当时的表面环境也发生了重大变化$^{[40]}$。
\subsubsection{形成之后}
地球的大气和海洋由火山活动和气体逸出形成$^{[41]}$。来自这些源的水蒸气凝结成海洋，并由小行星、原行星和彗星带来的水与冰进一步补充$^{[42]}$。按照一种模型，自地球形成之初起，地球上就可能已经有足以填满海洋的水量$^{[43]}$。在该模型中，大气中的温室气体阻止了当时只有当前 70\%光度的新生太阳导致海洋结冰$^{[44]}$。到距今约 35 亿年前，地球的磁场已经建立，这有助于防止大气被太阳风剥离$^{[45]}$。
\begin{figure}[ht]
\centering
\includegraphics[width=8cm]{./figures/e931dfc771eed8c4.png}
\caption{“淡橙色小点——对早期地球的一种印象画，呈现其富含甲烷并带有橙色色调的早期大气层$^{[46]}$。”} \label{fig_Earth_3}
\end{figure}
随着地球熔融的外层冷却，形成了第一批固体地壳，被认为成分为镁铁质$^{[47]}$。第一批大陆地壳成分更为长英质，是由这种镁铁质地壳部分熔融形成的$^{[47]}$。在太古宙沉积岩中发现的冥古宙时代锆石颗粒表明，至少部分长英质地壳早在 44 亿年前就已经存在，比地球形成仅晚约 1.4 亿年$^{[48]}$。

关于这一初始、体积很小的大陆地壳如何演化至如今的丰富状态，主要有两种模型$^{[49]}$：
(1) 相对稳定地持续增长直至今天$^{[50]}$，这一观点得到全球大陆地壳放射性定年数据的支持；
(2) 在太古宙期间大陆地壳体积迅速增长，形成了现今大陆地壳的主体$^{[51][52]}$，这一观点得到来自锆石中铪同位素和沉积岩中钕同位素证据的支持。通过大陆地壳的大规模循环再造，特别是在地球早期阶段，这两种模型及其证据是可以调和的$^{[53]}$。

新的大陆地壳是通过板块构造形成的，而板块构造最终由地球内部持续散失热量驱动。在数亿年的时间尺度上，构造力促使大陆地壳区域结合成超级大陆，随后又破裂分离。约在 7.5 亿年前，最早已知的超级大陆之一罗迪尼亚（Rodinia）开始裂解。之后大陆重新聚合，在 6 亿至 5.4 亿年前形成泛诺西亚（Pannotia），最终又形成盘古大陆（Pangaea），该超级大陆在约 1.8 亿年前开始裂解$^{[54]}$。

最近一次冰期循环模式始于约 4,000 万年前$^{[55]}$，并在约 300 万年前的更新世显著增强$^{[56]}$。自那以后，高纬度和中纬度地区经历了反复的冰川与间冰期循环，大约每 21,000 年、41,000 年和 100,000 年重复一次$^{[57]}$。最近一次冰期，俗称“最后冰期（Last Glacial Period）”，使大陆大部分区域至中纬度地区都被冰层覆盖，并于约 11,700 年前结束$^{[58]}$。
\subsubsection{生命起源与演化}
化学反应在约 40 亿年前产生了第一批能够自我复制的分子。约 5 亿年后，所有现存生命的最后共同祖先出现$^{[59]}$。光合作用的进化使生命形式能够直接利用太阳能。随之产生的分子氧（O(_2)）在大气中累积，并在与太阳紫外辐射的相互作用下，于高层大气中形成了保护性的臭氧层（O(_3)）$^{[60]}$。较小细胞被纳入较大细胞之中，产生了复杂的细胞——真核细胞$^{[61]}$。随着群体内细胞分化程度不断提高，真正的多细胞生物得以形成。在臭氧层吸收有害紫外辐射的帮助下，生命开始殖民地球表面$^{[62]}$。最早的生命化石证据包括：在西澳大利亚约 34.8 亿年前的砂岩中发现的微生物席化石$^{[63]}$；在格陵兰西部约 37 亿年前的变沉积岩中发现的生物成因石墨$^{[64]}$；在西澳大利亚约 41 亿年前岩石中发现的生物材料残留物$^{[65][66]}$。地球上最早的直接生命证据存在于约 34.5 亿年前的澳大利亚岩石中，其中显示了微生物的化石$^{[67][68]}$。
\begin{figure}[ht]
\centering
\includegraphics[width=8cm]{./figures/5a5b88da1cf2e867.png}
\caption{对太古宙（Archean）的艺术印象——这是地球形成之后的那个宙——画面中出现了圆形叠层石，它们是数十亿年前产生氧气的早期生命形式。晚期重轰炸（Late Heavy Bombardment）结束后，地壳已经冷却，富含水分的贫瘠地表上可以看到大陆和火山的痕迹，而当时的月球仍以今天一半的距离绕地球运行，看起来大约大 2.8 倍，并产生强烈的潮汐$^{[69]}$。} \label{fig_Earth_4}
\end{figure}
在新元古代（Neoproterozoic），约公元前 10 亿年至 5.39 亿年前（1000–539 Ma），地球表面可能有相当大部分被冰层覆盖。这个假说被称为“雪球地球”（Snowball Earth），它之所以备受关注，是因为这一时期发生在寒武纪大爆发之前，而寒武纪大爆发（约 5.35 亿年前，535 Ma）标志着多细胞生命复杂性显著增加的时期$^{[70][71]}$。

在寒武纪大爆发之后，地球上至少发生过五次大型物种大灭绝事件，以及众多较小的灭绝事件$^{[72]}$。除目前被提出的全新世灭绝（Holocene extinction event）之外，最近的一次发生在约 6600 万年前（66 Ma），当时一颗小行星撞击引发了非鸟类恐龙及其他大型爬行动物的灭绝，但小型动物——例如昆虫、哺乳动物、蜥蜴和鸟类——在很大程度上得以幸存。

在过去的 6600 万年（66 Mys）中，哺乳动物不断多样化；数百万年前，非洲的一种类人猿获得了直立行走的能力$^{[73][74]}$。直立行走使得工具使用成为可能，并促进交流，从而为更大的大脑提供所需的营养与刺激，进而推动了人类的演化。

农业的发展，随后是文明的出现，使人类开始深刻影响地球以及其他生命形式的性质与数量——这种影响一直持续至今$^{[75]}$。
\subsubsection{未来}
\begin{figure}[ht]
\centering
\includegraphics[width=8cm]{./figures/a4705ee315f9141e.png}
\caption{对大约 50–70 亿年后，当太阳进入红巨星阶段时灼热地球的构想图} \label{fig_Earth_5}
\end{figure}
地球预期的长期未来与太阳的未来紧密相关。在未来约 11 亿年内，太阳光度将增加 10\%，在未来约 35 亿年内将增加 40\%$^{[76]}$。地球表面温度的上升将加速无机碳循环，可能在约 1 亿至 9 亿年内将大气 CO(_2) 浓度降低至对现存植物致命的水平（对于 C4 光合作用约为 10 ppm）$^{[77][78]}$。植被的缺失将导致大气中氧气的消失，使现有动物生命无法存活$^{[79]}$。由于光度升高，地球的平均温度可能在 15 亿年后达到 100 ℃（212 ℉），所有海洋水分将会蒸发并逃逸到太空，从而可能在大约 16 亿至 30 亿年内触发失控温室效应$^{[80]}$。即便太阳保持稳定且永恒存在，由于随着地核缓慢冷却，中洋脊蒸汽喷发减少，现代海洋中的相当一部分水也将会沉入地幔$^{[80][81]}$。

太阳将在约 50 亿年后演化成红巨星。模型预测太阳将膨胀至大约 1 AU（1.5 亿千米；9,300 万英里），约为其当前半径的 250 倍$^{[76][82]}$。地球的命运尚不清楚。作为红巨星，太阳将失去大约 30\% 的质量，因此，如果没有潮汐效应，当太阳达到最大半径时，地球将移动到距离太阳 1.7 AU（2.5 亿千米；1.6 亿英里）的轨道；否则，在潮汐效应存在的情况下，它可能进入太阳大气并被汽化，其较重的元素沉入垂死太阳的核心$^{[76]}$。
\subsection{物理特性}
\subsubsection{大小与形状}
\begin{figure}[ht]
\centering
\includegraphics[width=8cm]{./figures/a3439116fb51e49f.png}
\caption{地球西半球，显示的是相对于地心的地形高度，而不是像常见地形图那样以平均海平面为参考。} \label{fig_Earth_6}
\end{figure}
地球由于处于流体静力平衡状态，呈现圆润的形状，平均直径为 12,742 公里（7,918 英里），是太阳系中第五大类地行星尺度天体，也是最大的类地天体。$^{[83][84]}$

由于地球自转，它呈现扁球体形状，在赤道发生隆起；其赤道直径比两极直径长 43 公里（27 英里）。$^{[85][86]}$ 地球的形状还存在局部地形变化；最大的局部变化，如马里亚纳海沟（低于当地海平面 10,925 米，约 35,843 英尺）会使地球平均半径缩短约 0.17\%，而珠穆朗玛峰（高于当地海平面 8,848 米，约 29,029 英尺）会使地球平均半径增加约 0.14\%。$^{[\text{n}6][89]}$ 由于地表距离地球质心在赤道隆起处最远，因此厄瓜多尔的钦博拉索火山顶（6,384.4 公里，或 3,967.1 英里）是地球表面距离地心最远的点。$^{[90][91]}$ 与较为刚性的陆地地形类似，海洋也呈现更具动态的地形变化。$^{[92]}$

为了测量地球局部地形的变化，大地测量学采用一种理想化的地球模型，得到大地水准面（geoid）形状。这种形状是通过将海洋理想化处理，即假设其完全覆盖地球且没有潮汐、风等扰动而得到的。其结果是一个光滑但不规则的大地水准面，提供一个平均海平面作为地形测量的参考高度。$^{[93]}$
\subsubsection{表面}
\begin{figure}[ht]
\centering
\includegraphics[width=8cm]{./figures/958849ebee4464c0.png}
\caption{地球的合成影像，清晰可分辨其不同类型的表面：海洋（蓝色）占据地球表面的主导区域；非洲大陆从葱郁（绿色）到干旱（棕色）的土地分布明显；而地球极地冰层以南极海冰（灰色）覆盖南大洋，以及南极冰盖（白色）覆盖南极大陆。} \label{fig_Earth_7}
\end{figure}
地球表面是大气层与固体地壳及海洋之间的边界。按此定义，其面积约为 5.1 亿平方千米（1.97 亿平方英里）$^{[11]}$。地球可按纬度分为北半球和南半球，也可按经度分为东半球和西半球。

地球表面的大部分为海洋水体：70.8\%，即 3.61 亿平方千米（1.39 亿平方英里）$^{[11]}$。这片巨大的咸水体通常称为世界海洋$^{[94][95]}$，使得拥有动态水圈的地球成为一个真正的水之星球$^{[96][97]}$，或海洋世界$^{[98][99]}$。事实上，在地球早期历史中，海洋可能曾完全覆盖地球$^{[100]}$。世界海洋通常被划分为太平洋、大西洋、印度洋、南冰洋和北冰洋从最大到最小依次排列。海洋覆盖着地球的大洋地壳，其中大陆架海在较小程度上覆盖着大陆地壳的大陆架。大洋地壳形成了大型的海洋盆地，具有诸如：深海平原,海底山,海底火山 $^{[85]}$,海沟,海底峡谷,大洋高原,以及贯穿全球的大洋中脊系统$^{[101]}$在地球极地区域，海洋表面覆盖着季节性变化的海冰，这些海冰常与极地陆地、永久冻土及冰盖相连，形成极地冰盖。
\begin{figure}[ht]
\centering
\includegraphics[width=8cm]{./figures/55e03c2e7e1d60d7.png}
\caption{地壳地形起伏} \label{fig_Earth_9}
\end{figure}
地球的陆地覆盖了地表的 29.2\%，约 $1.49\times10^{8},\text{km}^2$（5,800 万平方英里）。这些陆地包括全球范围内大量岛屿，但大多数陆地面积由四大洲块构成（按面积从大到小排列）：非洲－欧亚大陆、美洲、南极洲和澳大利亚 $^{[102][103][104]}$。这些大陆块进一步被划分并归类为各大洲。陆地表面的地形变化巨大，包括山区、沙漠、平原、高原以及其他地貌形态。陆地的海拔从最低的 −418 m（−1,371 英尺，位于死海）到最高的 8,848 m（29,029 英尺，珠穆朗玛峰峰顶）。陆地平均海拔约为 797 m（2,615 英尺）$^{[105]}$。

陆地表面可能被地表水、积雪、冰层、人造结构或植被覆盖。地球上大多数陆地具有植被覆盖 $^{[106]}$，但也有相当数量的陆地为冰盖（10\%，不包括同样面积巨大的永久冻土下的陆地 $^{[107][108]}$）或沙漠（33\% $^{[109]}$）。

土壤圈（pedosphere）是地球陆地表面最外层，由土壤构成，并受土壤形成过程影响。土壤对于陆地可耕性至关重要。地球可耕地约占陆地表面的 10.7\%，其中 1.3\% 为永久农田 $^{[110][111]}$。全球耕地面积估计约为 $1.67\times10^{7},\text{km}^2$（640 万平方英里），牧场面积约为 $3.35\times10^{7},\text{km}^2$（1,290 万平方英里）$^{[112]}$。

陆地表面与海底构成了地壳的顶部，并与上地幔部分共同形成岩石圈。地壳可分为洋壳和陆壳。洋壳在海底沉积物之下主要由玄武岩组成，而陆壳可能包含密度更低的岩石，如花岗岩、沉积岩和变质岩 $^{[113]}$。大陆表面约有 75\% 被沉积岩覆盖，尽管沉积岩仅占整个地壳质量的约 5\% $^{[114]}$。

地表地形包括海洋表面的地形及陆地形态。海底地形平均水深约 4 km，其地貌变化程度与海平面以上的陆地同样丰富。地球表面持续受到内部板块构造过程的塑造，包括地震与火山活动；以及由冰、水、风和温度驱动的风化与侵蚀；还包括生物过程，如生物量的生长与分解形成土壤 $^{[115][116]}$。
\subsubsection{板块构造}
\begin{figure}[ht]
\centering
\includegraphics[width=8cm]{./figures/c9553aaef0836668.png}
\caption{} \label{fig_Earth_10}
\end{figure}
地球地壳和上地幔构成的机械刚性外层，即岩石圈，被划分为若干构造板块。这些板块是彼此相对运动的刚性单元，并在三种类型的板块边界处发生移动：在聚合边界，两块板块相互靠拢；在张裂边界，两块板块被拉开；在转换边界，两块板块沿水平方向相互错动。在这些板块边界处，可能发生地震、火山活动、造山运动以及海沟的形成 $^{[118]}$。这些构造板块“漂浮”在软流圈之上，软流圈是上地幔中固态但黏度较低的部分，可以随板块一起流动和移动 $^{[119]}$。

随着构造板块的移动，洋壳在聚合边界处被俯冲到板块的前缘之下。同时，地幔物质在张裂边界处上涌，形成大洋中脊。这些过程共同作用，使洋壳不断回收进入地幔。由于这种循环，绝大部分海洋底部的年龄少于 1 亿年。年龄最老的洋壳位于西太平洋，估计约有 2 亿年 $^{[120][121]}$。相比之下，最早测定的陆壳年龄为 40.30 亿年 $^{[122]}$，虽然在始太古代沉积岩中保存的碎屑锆石显示年龄可达 44 亿年，表明当时至少已经存在少量大陆地壳 $^{[48]}$。

七大主要板块包括太平洋板块、北美板块、欧亚板块、非洲板块、南极板块、印－澳板块和南美板块。其他著名板块包括阿拉伯板块、加勒比板块、位于南美洲西海岸外侧的纳斯卡板块，以及南大西洋的斯科舍板块。澳大利亚板块在距今 5,000–5,500 万年前与印度板块融合。移动速度最快的是洋壳板块，例如科科斯板块以每年 75 mm（3.0 英寸/年）的速度前进 $^{[123]}$，太平洋板块以每年 52–69 mm（2.0–2.7 英寸/年）移动。相反，移动速度最慢的是南美板块，典型速度约为每年 10.6 mm（0.42 英寸/年）$^{[124]}$。
\subsubsection{内部结构}
\begin{figure}[ht]
\centering
\includegraphics[width=6cm]{./figures/ff301f5286e7d7b9.png}
\caption{} \label{fig_Earth_11}
\end{figure}
地球的内部结构与其他类地行星类似，根据化学或物理（流变学）性质被划分为不同层次。最外层是成分上独特的硅酸盐固态地壳，其下方是高黏度的固态地幔。地壳与地幔之间由莫霍不连续面（Moho面）分隔 $^{[127]}$。地壳厚度在海洋之下约为 6 公里（3.7 英里），而在大陆地区约为 30–50 公里（19–31 英里）。地壳与上地幔寒冷、坚硬的顶部共同构成岩石圈，岩石圈被分割为彼此独立运动的构造板块 $^{[128]}$。

在岩石圈之下是软流圈，它是相对低黏度的一层，岩石圈“漂浮”其上。在距地表约 410 公里和 660 公里处（250 与 410 英里），地幔内部会发生重要的晶体结构变化，形成了分隔上地幔与下地幔的过渡带。地幔以下是一层极低黏度的液态外核，其下是固态内核 $^{[129]}$。地球内核可能以略高于地球其余部分的角速度旋转，每年大约提前 0.1–0.5°，尽管也有提出更高或更低速率的模型 $^{[130]}$。内核半径约为地球半径的五分之一。密度随深度增加。在太阳系所有具行星尺度的天体中，地球是密度最高的对象。
\subsubsection{化学成分}
地球质量约为 $5.97\times 10^{24}kg$（5.970 Yg）。其主要成分为铁（按质量占 32.1\%）、氧（30.1\%）、硅（15.1\%）、镁（13.9\%）、硫（2.9\%）、镍（1.8\%）、钙（1.5\%）和铝（1.4\%），其余 1.2\% 为其他元素的微量成分。由于重力分异作用，地核主要由密度较高的元素组成：铁（88.8\%），并含少量镍（5.8\%）、硫（4.5\%）及不足 1\% 的微量元素 $^{[131][47]}$。

地壳中最常见的岩石成分是氧化物。超过 99\% 的地壳由十一种元素的各种氧化物组成，主要包括含硅的氧化物（硅酸盐矿物）以及含铝、铁、钙、镁、钾或钠的氧化物 $^{[132][131]}$。
\subsubsection{内部热量}
\begin{figure}[ht]
\centering
\includegraphics[width=8cm]{./figures/2d1016cbebea8e6a.png}
\caption{地球内部向地壳表面传热分布图，大部分热流集中在大洋中脊附近} \label{fig_Earth_12}
\end{figure}
地球内部热量的主要来源包括原始热（地球形成时遗留下来的热量）和放射性热（由放射性衰变产生的热量）$^{[133]}$。地球内部主要的产热同位素为钾-40、铀-238和钍-232$^{[134]}$。在地球中心，温度可能高达 6 000 °C（10 830 °F）$^{[135]}$，压力可达 360 GPa（5 200 万 psi）$^{[136]}$。由于大量热量来自放射性衰变，科学家推测，在地球早期历史中，在半衰期较短的同位素耗尽之前，地球的产热量要高得多。在约 30 亿年前，地球产生的热量可能是今天的两倍，这会加速地幔对流和板块构造，并促成如柯马提岩（komatiite）等如今罕见的火成岩的形成$^{[137][138]}$。

地球平均热流为 87 mW/m$^{2}$，全球总热流约为 $4.42\times 10^{13}$ W$^{[139]}$。部分来自地核的热能通过地幔柱输送至地壳，这是一种由高温岩石上涌构成的对流形式。地幔柱可以产生热点和洪流玄武岩$^{[140]}$。更多的热量则通过板块构造流失，即与海岭相关的地幔上涌过程。最后一种主要的失热方式是通过岩石圈传导，其中大部分发生在海洋下方$^{[141]}$。
\subsubsection{重力场}
参见主条目：Earth 的重力
地球重力是由于地球内部质量分布而赋予物体的加速度。在地球表面附近，重力加速度约为 $9.8 m/s^{2}$（$32 ft/s^{2}$）。地形、地质以及更深层构造的局部差异会导致地球重力场的局部及区域性变化，即所谓的重力异常$^{[142]}$。
\subsubsection{磁场}
\begin{figure}[ht]
\centering
\includegraphics[width=8cm]{./figures/4946ce6297772698.png}
\caption{地球磁层示意图，显示太阳风从左向右流动。} \label{fig_Earth_13}
\end{figure}
地球磁场的主要部分是在地核中产生的，那里发生着一种动力发电机过程，将由热和成分驱动的对流的动能转化为电能与磁场能。磁场从地核向外延伸，穿过地幔，一直延伸到地球表面，在那里大致呈偶极子形态。偶极子的两极位于地理南北极附近。在磁赤道，地表的磁场强度为$3.05\times10^{-5},\mathrm{T}$，在 2000 纪元时具有$7.79\times10^{22},\mathrm{A,m^2}$的磁偶极矩，并且每世纪大约减弱近 6\%（尽管仍然强于长期平均值）$^{[143]}$。地核中的对流运动是混沌的；磁极会漂移并周期性改变方向。这导致主磁场的长期变化，以及每隔数百万年平均发生数次的不规则磁场反转。最近一次反转发生在大约 70 万年前$^{[144][145]}$。

地球磁场在空间中的延伸范围定义了磁层。太阳风中的离子和电子会被磁层偏转；太阳风压力会压缩朝向太阳一侧的磁层至约 10 个地球半径，并把夜侧磁层延伸成一条长长的尾部$^{[146]}$。由于太阳风速度大于扰动在太阳风中传播的波速，在日侧磁层前方、太阳风内部会形成超音速弓形激波$^{[147]}$。

带电粒子被约束在磁层中；等离子层（plasmasphere）由低能粒子组成，它们基本上沿着磁力线随地球自转而运动$^{[148][149]}$；环电流（ring current）由相对于地磁场漂移的中能粒子组成，但其轨迹仍主要受磁场主导$^{[150]}$；范艾伦辐射带（Van Allen radiation belts）由高能粒子组成，其运动几乎是随机的，但仍被限制在磁层内部$^{[151][152]}$。在磁暴和次暴期间，带电粒子可能从外磁层，特别是磁尾区域偏转，沿磁力线被输送到地球电离层，在那里大气原子会被激发和电离，从而产生极光$^{[153]}$。
\subsection{轨道与自转}
\subsubsection{自转}
地球相对于太阳的自转周期——即其平均太阳日——为平均太阳时的 86,400 秒（86,400.0025 SI 秒）$^{[154]}$。由于潮汐减速作用，地球的太阳日现在略长于 19 世纪，因此每天的长度比平均太阳日长 0 到 2 毫秒不等$^{[155][156]}$。

地球相对于恒星（即固定恒星）的自转周期，由国际地球自转与参考系统服务（IERS）称为恒星日，为平均太阳时（UT1）的 86,164.0989 秒，即 $23^h56^m 4.0989^s$,$^{[2][n,10]}$。地球相对于岁差或移动的平均春分点（当太阳位于赤道 90° 时）的自转周期为平均太阳时（UT1）的 86,164.0905 秒（$23^h56^m 4.0905^s$）$^{[2]}$。因此，恒星日比岁差日短约 8.4 ms$^{[157]}$。

除了大气层内的流星和低轨卫星以外，地球天空中天体的主要视运动方向为向西，视运动速率为 15°/h = 15′/min。对于靠近天球赤道的天体而言，这相当于每两分钟一个太阳或月亮的视直径；从地球表面观察，太阳和月亮的视大小大致相同$^{[158][159]}$。
\subsubsection{轨道}
\begin{figure}[ht]
\centering
\includegraphics[width=8cm]{./figures/5ba001e8ae217704.png}
\caption{夸张示意图显示了地球绕太阳的椭圆轨道，并标明轨道极值点（远日点与近日点）与四个季节的极值点——春分、秋分与夏至、冬至——并不相同。} \label{fig_Earth_14}
\end{figure}
地球围绕太阳公转，使地球成为距离太阳第三近的行星，并属于太阳系内侧行星。地球的平均轨道距离约为 1.5 亿千米（9,300 万英里），这是天文单位（AU）的基准，约等于 8.3 分钟的光行程，或约为地月距离的 380 倍。地球每 365.2564 个平均太阳日绕太阳运行一周，即一个恒星年。由于太阳在地球天空中以每天约 1° 向东的视运动——相当于每 12 小时一个太阳或月亮的视直径——平均而言，地球需要 24 小时（一个太阳日）才能完成自转，使太阳重新返回子午线。

地球的平均轨道速度为 29.7827 千米/秒（107,218 千米/小时；66,622 英里/小时），这一速度足以在 7 分钟内走完地球直径约 12,742 千米（7,918 英里），并在约 3.5 小时内走完地月距离 384,400 千米（238,855 英里）$^{[3]}$。

月球与地球每 27.32 天围绕共同的质心相对于背景恒星运行一周。当与地月系统绕太阳的共同轨道相结合，从朔月到朔月的会合月周期为 29.53 天。从天球北极观测，地球、月球及其自转均呈逆时针方向运动。从太阳和地球北极上方的观测点看，地球绕太阳亦呈逆时针方向公转。轨道平面与自转轴面并非完全对齐：地球自转轴相对于地日平面（黄道面）的垂直方向倾斜约 $23.44$，而地月轨道平面相对于地日平面倾斜最多可达 $\pm5.1^\circ$。若无此倾角，每隔两周就会发生一次食现象，在月食与日食之间交替$^{[3][160]}$。

地球的希尔球（Hill sphere），亦即地球的引力影响范围半径，约为 150 万千米（93 万英里）$^{[161][n,11]}$。这是地球引力相对于更远处太阳和行星引力仍为主导的最大距离。天体必须在此半径内绕地球运行，否则可能因太阳的引力摄动而脱离轨道$^{[161]}$。地球连同整个太阳系位于银河系之中，并以距银河系中心约 28,000 光年的距离绕其运行。地球位于猎户臂之中，约在银河盘面上方 20 光年的位置$^{[162]}$。
\subsubsection{自转轴倾角与季节}


